\documentclass[11pt]{report}
\usepackage[margin=1in]{geometry}
\usepackage{graphicx,siunitx,booktabs,tabularx,array,amsmath,mathabx,natbib,url}
\graphicspath{{D:/Github/Bipedal_Robot/Documentation/Reports and Papers/Dissertation/Notes/advisor_revision_20261023/REVIEW_INBOX/REVIEWED_20261026/figures/}}
\makeatletter\newlabel{app:muscle_calcs}{{X}{X}}\makeatother
\makeatletter\newlabel{ch:muscles}{{X}{X}}\makeatother
\makeatletter\newlabel{fig:PulleyMargins_resolved}{{X}{X}}\makeatother
\makeatletter\newlabel{fig:PulleySurfaces}{{X}{X}}\makeatother
\makeatletter\newlabel{sec:app_routing_extensions}{{X}{X}}\makeatother
\makeatletter\newlabel{sec:sns_existing_tools}{{X}{X}}\makeatother
\begin{document}
\chapter*{Targeted revision fragments: review proof}
This proof checks proposed fragments and figure placement. It is not the assembled dissertation. Cross-references to material outside the packet are represented by the placeholder X.
Researchers studying neuromechanical control need a reusable way to assemble neural controllers with human musculoskeletal and artificial-muscle models and evaluate their behavior under specified support, stimulation, and perturbation conditions. I addressed this need by developing the following models and tools from the existing simulation packages:
\begin{enumerate}
    \item A custom Li--Deng walker and controller, developed in AnimatLab and implemented in MuJoCo with SNS-Toolbox, with matched implementations of its contact-driven neural core in NumPy and Simulink.
    \item Gait2392-based neuromechanical models with alternative neural controllers, and a robot counterpart with BPA routing and the actuator-force model from Chapter~\ref{ch:muscles}.
    \item The SNS-Simscape block library and GUI for block programming neural pathways, including sensory, neuron, synapse, activation, and custom BPA-force blocks, together with coupled reflex and CPG demonstrations.
    \item A SolidWorks-derived robot assembly in Simscape Multibody coupled to a neural controller and BPA actuators using the force model from Chapter~\ref{ch:muscles}.
    \item Integration of SNS-Simscape controllers with MuJoCo through the existing MathWorks bidirectional bridge, and model-import tools for bringing musculoskeletal and CAD-derived mechanisms into the MATLAB environment.
    \item An automated curriculum and experiment workflow that combines staged parameter searches with adjustable support, neural stimuli, pathway interventions, mechanical perturbations, and recorded neural and mechanical outputs.
    \item Implementations of literature-derived circuits and selected spiking-neuron counterparts for comparing neural mechanisms and coding assumptions.
\end{enumerate}
OpenSim, MuJoCo, SNS-Toolbox, Optuna, Simulink, and the MathWorks bridge provide the underlying software; the contributions here are the model implementations, custom blocks, integrations, and evaluation procedures. Table~\ref{tab:sns_capability} presents selected functional-verification tests and proof-of-concept demonstrations illustrating the capabilities of this toolset.

\clearpage
\begin{table}[htbp]
\caption[Tools developed and representative demonstrations]{Tools developed in this project and representative functional-verification tests and proof-of-concept demonstrations. The examples illustrate the toolset rather than exhaustively evaluating every implementation; each result refers to the named circuit or mechanical model.}
\label{tab:sns_capability}
\small
\begin{tabularx}{\textwidth}{@{}>{\raggedright\arraybackslash}p{0.23\textwidth}*{2}{>{\raggedright\arraybackslash}X}@{}}
\toprule
\textbf{Tool developed} & \textbf{Function} & \textbf{Representative demonstration} \\
\midrule
Li--Deng/W2L model implementations & Combine the walker mechanics and neural architecture; transfer selected controller components between environments & Matched contact-driven rhythm cores reproduce \qty{1.000}{\second} periods with trace correlation at least 0.99996 \\
Gait2392/SNS model implementations & Evaluate alternative controllers on the human reference and custom actuator/routing models & Individual-actuator tests check torque directions; selected human-model controllers produce bilateral stepping-like activity under support \\
SNS-Simscape library and GUI & Construct neural pathways graphically and connect sensory, neural, and actuator blocks & The loaded-arm reflex demo gives \ang{2.00} sag with feedback and \ang{7.85} without it; actuator blocks agree with reference force calculations \\
W2L BPA-force integration & Apply the Festo force model within the coupled simulation & The 12-actuator BPA model demonstrates air stepping and retained-support standing, with peak pressure \qty{372}{\kPa} \\
SNS-Simscape/MuJoCo integration & Exchange activation commands and mechanical measurements through the MathWorks bridge & Single-muscle and control-rate tests reproduce the Python-reference sensor readback over the scored windows \\
Curriculum and experiment drivers & Automate staged tuning, support changes, neural stimuli, and mechanical perturbations & Saved stage-specific trials, candidate replays, and fixed-controller pathway and body-change comparisons demonstrate the workflow \\
Literature-circuit and spiking implementations & Test selected published neural behaviors and compare graded and spiking formulations & The reduced Rybak circuit passes 13 drive/deletion checks; selected spiking counterparts have paired implementation tests \\
\bottomrule
\end{tabularx}
\end{table}

\clearpage
My first conversion attempt used the robot model, which combined BPA routing on the Gait2392 skeleton with unchanged human muscle parameters. MyoConverter's muscle-equilibrium checks rejected infeasible combinations of routed length and muscle parameters. I therefore separated conversion of the human reference from reconstruction of the robot routing. This distinction matters because an error caused by an artificial-actuator route should not be interpreted as a property of the original human model.

The human-model conversion also required specific repairs. Three small hip rotators per side, the quadratus femoris, gemelli, and piriformis, were disabled at model load. The converted model retains 92 actuator entries, but these six contribute no force. The model retains the OpenSim and International Society of Biomechanics coordinate convention; frame transformations and individual-actuator tests check that joint motion and torque directions have the intended signs.

The source represents patella motion through muscle insertion locations that vary with knee angle. In the converted model, limits on the joints representing these moving path points conflicted with their coupling constraints. Removing those conflicting limits allowed the coupled path points to follow their prescribed trajectories. The rectus femoris path also required a patella-tracking via point to restore its knee-extension action. Explicit collision pairs introduced unwanted contacts, so contact was restricted to the foot--ground pairs. The anatomical joint limits also had to be re-enabled; these driver joints were treated separately from the coupled path-point joints. Passive joint torques supplemented the joint behavior. These repairs explain why conversion needed checks beyond successful loading: a model could load while producing an incorrect torque direction or unintended contact force.

Comparisons with the OpenSim attachment paths and single-actuator tests checked the transferred geometry and force directions. They do not establish that every muscle reproduces the source force throughout its range or that the controller produces walking. The neural and mechanical models exchange signals at a \qty{2}{\milli\second} solver time step. Appendix~\ref{sec:app_conversion_repairs} records the individual repairs and their scope.


\clearpage
\section{OpenSim-to-MuJoCo Conversion Repairs}\label{sec:app_conversion_repairs}

The conversion work in Section~\ref{sec:sns_existing_tools} involved two distinct problems: importing the robot's artificial-muscle routing with human muscle parameters, and repairing the converted human reference used for controller development. The first combination failed muscle-equilibrium checks because the routed lengths were incompatible with the retained human parameters. The subsequent human-model repairs described below belong to that reference-model loading path; they are not evidence that the BPA-routed model reproduces every human muscle force.

\paragraph{Joint directions and coordinate conventions.}
Single-actuator tests exposed reversed hip-flexion and hip-adduction actions. The corresponding hinge-axis directions were negated so that the muscle actions agreed with the adopted OpenSim conventions. The knee required a different treatment: its converted flexion-negative convention was retained. Rotating the knee without updating its coupled translational coordinates gives a misleading geometry comparison, so those coordinates must follow the prescribed coupling during a pose or actuator test. These are specific joint and path checks, rather than a change from a left-handed to a right-handed model.

\paragraph{Small muscles and force scales.}
The quadratus femoris, gemelli, and piriformis actuators were disabled on both sides. Their six entries remain in the 92-actuator model, with no force contribution. A separate check found eight muscle-actuator entries with an active maximum-force scale of only \qty{1}{\newton}. Their active-force scales were restored from the source-model values. This correction concerns the human muscle actuators, not the BPA force law.

\paragraph{Patella-related paths and coupled joints.}
The converted moving muscle-path points use joints coupled to knee angle. Their sampled limits conflicted with the coupling constraints away from the sampled pose, producing large constraint forces and restricted motion. The coupled path-point joints were therefore left without those conflicting limits while retaining their prescribed trajectories. The anatomical driver joints were treated separately: limits were re-enabled on 14 hinges using the converted ranges. The rectus femoris route also lacked the patella tracking needed for knee-extension action; its intermediate path point was replaced by the vastus medialis patella-tracking point. The resulting force comparison and residual routing differences are discussed in Appendix~\ref{app:muscle_calcs}.

\paragraph{Contact and virtual geometry.}
The converted contact section contained 19 explicit ground--body pairs. These predefined pairs are independent of the individual geometry contact flags, so changing those flags alone did not remove ground contact in a suspended test. Suspended tests remove the explicit contact pairs; ground-contact tests retain only the six foot--ground pairs. Virtual muscle-path marker geometry was also prevented from colliding. This separates intended foot loading from unintended contact involving path markers or other body segments.

The repairs establish specific geometry, sign, constraint, and contact corrections. They do not remove the rigid-tendon approximation or establish complete equivalence with the OpenSim model. The appropriate result remains the corresponding test: correct muscle action, a valid coupled trajectory, intended ground contact, or agreement in the measured force comparison. None alone establishes walking performance.


\clearpage
\subsection{Controlled Perturbation and Pathway Tests}\label{sec:w2l_perturbation}
The experiment drivers apply neural stimuli, external forces, pathway deletions, and body-mass changes to a fixed controller configuration. Replaying all stored controller fields gives 17 left and 16 right strides in the control condition. The corresponding counts are 16/16 after crossed-heel deletion, 13/13 after Ia-chain deletion, 16/14 after Renshaw-cell deletion, 17/15 after heel-cascade deletion, and 13/12 after commissural deletion. Removing Ib load feedback stops stepping in both legs. Thus, Ib feedback is necessary for stepping at this tested operating point; the other deletions retain stepping with varying counts.

The neural stimuli and external-force tests also retain bilateral stepping. Counts alone do not establish balance recovery: the stronger stimulus and both push conditions exceed the \ang{45} tilt criterion. Across ten body-mass conditions, stride counts span 32--38 per leg and the lowest recorded pelvis height is \qty{0.8981}{\meter}. Nine conditions exceed the tilt criterion; one meets the full survival gate. These comparisons demonstrate controlled changes in neural and mechanical conditions while distinguishing continued stepping from postural stability.

\begin{figure}[p]
\centering
\includegraphics[width=\textwidth]{w2l_perturbation_fully_applied.pdf}
\caption[Controlled W2L perturbation and pathway tests]{Supported W2L tests with all stored controller fields applied. The control produces 17/16 left/right strides. Ib deletion stops stepping; commissural deletion retains 13/12 strides. Stimulus and push trials retain stepping, although several exceed the tilt criterion. The controller replay does not constitute a new optimization with the corrected parameter forwarding.}
\label{fig:w2l_perturbation}
\end{figure}

\clearpage
At fixed controller and drive settings, removal of contact and Ib force-feedback pathways reduced stepping in the supported W2L model from 33 left and 28 right cycles to one left and no right cycles. Forward displacement decreased from \qty{1.16}{\meter} to \qty{0.09}{\meter}. These measures cover the \qty{25.499}{\second} scored window after settling within each \qty{30}{\second} simulation. The harness configuration was unchanged, although its measured support contribution differed between conditions. Ia and II pathways remained active. This comparison demonstrates the contribution of the removed pathways at this operating point; it uses the standard muscle-actuated W2L model.
\begin{figure}[t]
\centering
\includegraphics[width=0.85\textwidth]{ablation_pair_metrics_20261026.pdf}
\caption[Contact and force-feedback ablation]{Contact and Ib force-feedback ablation at fixed controller and drive settings in the supported W2L model. Counts and displacement refer to the post-settling scored window. The retained Ia and II pathways distinguish this test from complete deafferentation.}
\label{fig:sns_ablation_pair}
\end{figure}

\clearpage
The archived rat-hindlimb model of \citet{deng_neuromechanical_2019} provides an original-engine demonstration of the experimental workflow. The model was exported to standalone AnimatLab simulation form and instrumented with neural-voltage charts. In a \qty{10.1}{\second} run without external neural drive, rhythm-generator extensor and flexor populations alternate with an approximately \qty{0.23}{\second} period, while the two legs fire in phase. The baseline retains the archived model's mechanical torque input.

A \qty{2}{\nano\ampere} current was applied to the left rhythm-generator flexor from \SIrange{2}{2.1}{\second}, using the stimulation timing reported by \citet{deng_biarticular_2023}. A flexor burst begins at \qty{2.000}{\second}, compared with \qty{2.433}{\second} in the baseline. The post-stimulus period is \qty{0.224}{\second}, compared with \qty{0.221}{\second} in the baseline. Applying the same current and timing to the knee pattern-formation flexor produces a post-stimulus rhythm-generator phase difference of \qty{-2.3 \pm 2.4}{\percent} of the baseline period. An additional, longer knee-pattern-formation stimulus from \SIrange{2}{2.5}{\second} illustrates the motor-output response in Figure~\ref{fig:deng_original_perturbation}D. These are observations from the archived antagonist model in its original engine, rather than a quantitative validation of the Simulink translation or a reproduction of the later biarticular model's results.
\begin{figure}[p]
\centering
\includegraphics[width=\textwidth]{deng_original_reviewed.pdf}
\caption[Archived Deng-model neural perturbations]{Neural perturbations in the archived antagonist model, simulated in its original AnimatLab engine. Solid lines show the baseline; dashed lines show stimulated conditions. (A, B)~Left rhythm-generator flexor and extensor responses to rhythm-generator flexor stimulation. (C)~Left rhythm-generator extensor response to knee pattern-formation flexor stimulation. The current is \qty{2}{\nano\ampere} from \SIrange{2}{2.1}{\second} in A--C. (D)~Left knee extensor motoneuron response to an extended knee pattern-formation flexor stimulus from \SIrange{2}{2.5}{\second}. Shading marks each panel's actual stimulus window.}
\label{fig:deng_original_perturbation}
\end{figure}

\clearpage
% P04: replaces the whole body of
% \subsection{Mono- and Bi-Pulley Tendon Routing}\label{sec:disc_pulley}
\subsection{Mono- and Bi-Pulley Tendon Routing}\label{sec:disc_pulley}

The pulley algorithms extend the route search space by allowing the artificial tendon to be reeved over pulleys before insertion, and the naming refers to how many degrees of freedom dominate the routing calculation: the mono-pulley search computes the route on the knee's dominant degree of freedom, and the bi-pulley search computes it on the two dominant degrees of freedom of the biarticular evaluation planes. The purpose of the reversed pulley is a trade of force for range of motion: the two-part reversed tackle halves the insertion-side force while doubling the insertion travel, which can extend the angle range over which a route meets its torque requirement. In the knee-flexor route the pulley acts as a redirect in this configuration rather than a force multiplier, and the two optimized routing configurations differ by more than the pulley alone because each was solved under its own objective: peak predicted knee-axis torque moved from \qty{81.23}{\newton\meter} to \qty{85.74}{\newton\meter}, with off-axis torque held to \qty{0.803}{\newton\meter}. Appendix~\ref{sec:app_routing_extensions} records the optimizer diagnostics and the broader route searches.

A subsequent transmission study re-solved actuator sizes for each tackle configuration at five anatomical muscle sites, making the force-for-travel trade explicit in the sizing. The corresponding torque and margin surfaces are shown in Figures~\ref{fig:PulleySurfaces} and~\ref{fig:PulleyMargins_resolved}. These calculations use the moving-force-point routing, with the adjacent joint fixed for the biarticular evaluation planes; they do not establish simultaneous performance across both joints. Physical verification of the pulley configurations remains future work.

\clearpage
\begin{figure}[p]
\centering
\makebox[\textwidth][l]{\textbf{(A)}}\par
\includegraphics[width=\textwidth]{Ext10mm_52cm-eps-converted-to.pdf}\par\medskip
\makebox[\textwidth][l]{\textbf{(B)}}\par
\includegraphics[width=0.52\textwidth]{KneeExt_10mm_passthru.png}
\caption{Measured versus expected isometric torque using a $\diameter$\qty{10}{mm}, $l_{0}=\qty{51.8}{\cm}$ extensor BPA (A), and the corresponding routing configuration (B). In (A), blue dots show measured results, the dashed black line is the original prediction, the solid magenta line is the optimized prediction, and yellow dots show the hybrid calculation. In (B), the BPA passes between the femoral condyles close to the instantaneous center of rotation.}
\label{fig:Ext10mm_52cm}\label{fig:Ext_passthru}
\end{figure}

\clearpage
\begin{figure}[p]
\centering
\refstepcounter{figure}\label{fig:extensor_physical_final}
\includegraphics[width=0.97\textwidth]{extensor_20mm_physical.pdf}
\end{figure}
\begin{figure}[p]
\centering
\refstepcounter{figure}\label{fig:flexor_physical_final}
\includegraphics[width=0.97\textwidth]{flexor_20mm_physical.pdf}
\end{figure}

\clearpage
\bibliographystyle{plainnat}
\bibliography{proof_refs}
\end{document}
